\chapter{System Identification and Calibration}
\label{ch:calibration}

% TODO: this is the thesis's strongest evidentiary chapter - consider
% pulling the six-point sweep and the readMode before/after table into
% actual plotted figures (velocity-vs-time traces) rather than tables
% only, if you still have the raw CSV/rosbag data from the sessions
% referenced here.

This chapter documents the longest and most consequential
investigation of the project. A real-hardware symptom emerged in
which the vehicle, commanded to travel a fixed distance and then
stop, consistently travelled substantially further before actually
coming to rest, in the worst observed case roughly three times as far.
What began as a search for a physical explanation (drivetrain
coasting) ended in the discovery of a stale-data artefact in the
vendor hardware abstraction layer described in
Section~\ref{sec:qcar-platform}. The true physical residual turned out
to be small once that artefact was removed. The chapter also
documents a control-design attempt that was tested on physical
hardware without adequate prior validation and produced a genuine,
dangerous oscillation. It is included here in full because it is as
methodologically important as the investigation's eventual successful
resolution.

\section{Problem Statement}
\label{sec:stop-distance-problem}

With the bridge of Chapter~\ref{ch:hardware-bridge} operational, a
simple closed-loop test was used throughout this investigation: drive
the vehicle at a fixed commanded speed toward a target distance,
using odometry to detect when the target is crossed and stop.
Repeated trials showed the vehicle travelling substantially further
than the commanded target before settling, with the overshoot
growing, not shrinking, across repeated sessions under otherwise
similar conditions. This was an early indication that the effect was
systematic and reproducible rather than simple measurement noise.

\section{Ruling Out Simpler Explanations}
\label{sec:ruled-out}

Several plausible causes were tested directly and eliminated before
the investigation's eventual resolution:

\begin{itemize}
  \item \textbf{AMCL/localisation.} The original suspect. Investigating
    it produced one genuine, separately valid fix. This raised AMCL's
    \texttt{recovery\_alpha\_fast}/\texttt{recovery\_alpha\_slow} from
    \texttt{0.0} to \texttt{0.1}/\texttt{0.001}, since a value of
    \texttt{0.0} allows the particle filter to freeze permanently once
    it (incorrectly) converges to zero uncertainty under real sensor
    noise. This failure mode is not exposed by a clean simulation
    environment. This fix was kept, but it does not explain the
    stop-distance symptom, which persisted after it was applied.
  \item \textbf{Command-delivery delay.} Measured directly on the
    QCar's own local clock (eliminating any network-latency
    contribution) at approximately \SI{0.2}{\second} from decision to
    applied throttle. This is normal and not large enough to explain
    metre-scale overshoot at the tested speeds.
  \item \textbf{A backlog of stale \texttt{/cmd\_vel} commands.}
    Hypothesised as queuing behind a blocking socket \texttt{send}
    call in the relay's single-threaded executor. Two increasingly
    careful fixes were built and deployed: a staleness timestamp
    check, then a full non-blocking sender-thread architecture. Both
    were tested directly on hardware. \textbf{Neither made any
    measurable difference}, ruling this out as the cause.
  \item \textbf{Logging/print buffering artefacts.} The timing of
    QCar-local debug samples during the ``coast'' period was a steady
    \SIrange{20}{35}{\milli\second} throughout, not a rapid burst
    consistent with a draining backlog -- ruled out.
\end{itemize}

\section{Active Braking as a Control-Design Attempt and a Safety Incident}
\label{sec:active-braking}

With the above ruled out, the leading hypothesis shifted to genuine
drivetrain (motor and gearbox) rotational inertia keeping the wheels
turning for a period after commanded throttle reached zero. This was
a physically plausible mechanism and one apparently consistent with
QCar-local encoder readings that showed real velocity persisting for
several seconds after a stop was commanded.

A corrective control law was implemented directly in
\texttt{qcar\_bridge.py}: a reverse-throttle brake pulse proportional
to the residual measured velocity, intended to actively cancel the
coasting motion. \textbf{This was deployed directly to physical
hardware without first being traced through by hand or tested in
simulation.} On its first test, it produced a genuine, dangerous
oscillation: the vehicle lurched forward and backward repeatedly,
reaching \SI{-0.88}{\meter\per\second} in reverse and still
accelerating. It had to be stopped manually via an SSH-issued kill
command mid-test. The root cause, in hindsight, is a textbook control
failure mode. A continuous proportional reaction to \emph{any}
nonzero residual velocity, including reverse velocity, with no
deadband and no explicit disarm condition, means every correction's
own overshoot looks to the controller like a fresh error requiring
correction. The result is an undamped hunting loop with no
mechanism to settle.

Since both attempts were reverted, no implementation of this control
law survives in the repository to quote directly. Listing~\ref{lst:active-braking-comment}
instead reproduces the comment left in its place in
\texttt{qcar\_bridge.py}, at the site the code once occupied. This is
the project's own permanent record of what was tried, why it failed
and why it must not be blindly re-attempted.

\begin{lstlisting}[language=Python,caption={Comment documenting the rejected active-braking design, left in place of the reverted code (\texttt{qcar\_bridge.py}).},label={lst:active-braking-comment}]
# Tried active braking (brief reverse-throttle pulse proportional to
# residual speed when a stop is commanded) to counter drivetrain coast -
# REJECTED, not just untuned. It's a proportional loop on velocity sign
# with no deadband and one-cycle-stale (~20ms) feedback: if the pulse
# overshoots past zero into real reverse motion, the next cycle reads a
# negative v and "corrects" with a forward push, which can itself
# overshoot - a genuine undamped hunting oscillation, not a bad gain
# value. Confirmed on hardware (twice - it was reverted once already
# before this comment, then re-tried and reproduced the same result):
# car visibly lurched forward/back repeatedly and fast, had to be
# killed manually. Do not re-add this without a real deadband/
# hysteresis or a one-shot (non-continuous) brake pulse design -
# retuning BRAKE_GAIN alone will not fix it.
\end{lstlisting}

\paragraph{A minimal model of why this design cannot be stable.}
The instability the comment above describes can be shown formally
with a deliberately simplified discrete-time model that keeps only
the mechanism responsible, not the full nonlinear vehicle dynamics.
Let $v_k$ be the residual measured velocity at control cycle $k$.
Let the brake law command a throttle proportional to that residual
with gain $K$: $u_k = -K v_k$. If the vehicle's velocity responds to
one cycle's throttle command with an effective per-cycle gain $a$
(lumping motor response, control period and vehicle mass into a
single constant and ignoring natural friction/drag for this minimal
model), the resulting closed-loop update is

\begin{equation}
  v_{k+1} = v_k + a\,u_k = v_k(1 - aK).
  \label{eq:brake-instability}
\end{equation}

This is a first-order linear recurrence. Its behaviour is
governed entirely by the magnitude of $(1 - aK)$. The residual
velocity decays toward zero only if $|1 - aK| < 1$, i.e.\
$0 < aK < 2$; it oscillates in sign every cycle if $aK > 1$.
Critically, that oscillation \emph{grows} in magnitude, without
bound, whenever $aK > 2$, since each successive $v_k$ is then
multiplied by a factor more negative than $-1$. This single
inequality is the formal statement of exactly what the source comment
describes qualitatively: a correction that overshoots past zero into
real reverse motion, prompting a next-cycle correction in the
opposite direction that itself overshoots by more. It also shows
directly why \emph{no} amount of retuning $K$ alone guarantees safety
here. $a$, the vehicle's true, real per-cycle response to a given
throttle command, was never independently measured before this
control law was deployed, so there was no way to have verified
$aK < 2$ in advance from the gain value alone.

A one-shot design avoids re-arming after the first correction, but
Equation~\ref{eq:brake-instability}'s single-step overshoot condition
does not require re-arming to manifest. If the very first correction
alone satisfies $aK > 2$, it already produces a reverse velocity
larger in magnitude than the original residual it was meant to
cancel. This is consistent with the one-shot retest reproducing the
same oscillation, rather than being protected by its disarm
condition. It is precisely the kind
of hand-traceable instability check Section~\ref{sec:methodology-overview}'s
standing rule asks for \emph{before} hardware deployment, not after.

A revised, ``one-shot'' design was implemented: braking fires exactly
once per drive-to-stop transition and permanently disarms the
instant measured velocity crosses zero, rather than re-arming from
residual drift. It \textbf{reproduced the same oscillation on
retest}. Given the repeated safety failure, this control approach was
abandoned entirely rather than iterated on further. Both modified
files were reverted to their exact pre-investigation state, verified
byte-identical to the known-good version by hash comparison. The
failure was documented directly in a source-code comment so that it
is not blindly re-attempted in a future session.

This incident is the clearest instance in this project of the general
principle stated in Section~\ref{sec:methodology-overview}. A
reactive feedback controller acting on a live sensor reading, on a
system already confirmed to have real actuator lag and inertia,
should be traced through by hand for an overshoot scenario or tested
in simulation first, before it is ever sent to physical hardware.
``The control theory seems sound'' was not, on its own,
sufficient justification here. This project treats that as a
standing rule for any future reactive control logic on this
platform, not merely a lesson learned in retrospect.

\section{Sim-versus-Real Comparison}
\label{sec:sim-vs-real}

Following a methodology suggested by my academic
supervisor, the identical closed-loop stop-distance test (drive $N$
metres at a fixed speed, stop the instant odometry crosses the
target) was reproduced in both the physical robot and the Gazebo
simulation of Chapter~\ref{ch:sim-stack}, using the same launch stack
with only the drive target switched from simulated to real hardware.

A six-point sweep across three target distances
(\SIlist{1.0;1.5;2.0}{\meter}) and three commanded speeds
(\SIlist{0.20;0.25;0.30}{\meter\per\second}) found real
hardware overshoot in the range \SIrange{0.58}{2.12}{\meter} across
two sessions roughly one hour apart. \emph{Every} tested condition
was worse in the second session, by a broadly consistent margin. This
is evidence of real, reproducible drift, plausibly battery voltage or
floor/tyre condition, rather than measurement noise.

The matched simulation sweep, over the same six conditions, showed
overshoot of only \SIrange{0.044}{0.067}{\meter} --
\SIrange{10}{30}{}$\times$ smaller and only weakly sensitive to
speed. This is in contrast to the real hardware overshoot, which
scaled roughly with the cube of speed. A dedicated, isolated
comparison at \SI{2}{\meter} / \SI{0.3}{\meter\per\second} sharpened
this further. The simulated vehicle settled at \SI{2.067}{\meter}
(an extra \SI{0.067}{\meter}), while the real vehicle settled at
\SI{3.91}{\meter} (an extra \SI{1.90}{\meter}). This is a
$\sim$28$\times$ difference for an identical software command. This
result strongly supports a real-hardware-specific controller/dynamics
cause, rather than something inherent to the stop-command approach
itself. It directly motivated the calibration work in the remainder
of this chapter.

\section{Throttle Calibration}
\label{sec:throttle-calibration}

\subsection{The Uncalibrated Gain Hypothesis}
\label{sec:throttle-gain-hypothesis}

\texttt{qcar\_bridge.py}'s mapping from \texttt{cmd\_vel}'s
\texttt{linear.x} (in \si{\meter\per\second}) to the HAL's PWM duty
cycle had, since early in the project, used a placeholder gain
(\texttt{THROTTLE\_GAIN = 1/3}), explicitly documented in its own
source comment as uncalibrated. Direct measurement from the sweep
data above confirmed this mattered. Commanding
\SI{0.3}{\meter\per\second} produced a real, QCar-local
encoder-measured cruise velocity averaging
$\approx$\SI{0.8}{\meter\per\second}, roughly $2.6$--$2.7\times$ the
commanded speed, consistent across two independent measurements.
Since coasting distance from momentum scales with the square of
velocity, this alone predicts roughly a $6.9\times$ larger coast
distance than a genuine \SI{0.3}{\meter\per\second} stop would
produce. This is a substantial fraction of the sim-vs-real gap found
above and a reminder that the sim-vs-real comparison in
Section~\ref{sec:sim-vs-real} was not, at that point, comparing truly
matched real speeds.

\subsection{A First Calibration Attempt That Discovered a Deadband, Then Reverted}
\label{sec:deadband-discovery}

A dedicated measurement tool, \texttt{check\_stop\_latency.py}, was
built to measure cruise velocity and stop latency in a single run. An
initial correction lowered \texttt{THROTTLE\_GAIN} by dividing by
the previously measured $2.67\times$ ratio. Deployed to hardware, it
immediately surfaced a second, independent, hard non-linearity. At
the old gain, a commanded \SI{0.25}{\meter\per\second} ($\sim$8\%
duty) sat completely motionless for $\approx$\SI{2.1}{\second}
before suddenly starting to move. This delay was later traced
(Section~\ref{sec:readmode-bug}) to a HAL buffering artefact rather
than a friction effect. \SIlist{0.20;0.15}{\meter\per\second}
($\sim$6.7\%/5\% duty) never moved at all across
\SIrange{8}{10}{\second} of commanding.

This second behaviour is zero motion regardless of how long the duty
is held. It is a genuine, independent physical deadband property of
the vehicle, not a linear gain error. It interacts badly with a
naive linear-gain correction: the lowered gain restored accurate
cruise-speed tracking at higher commanded speeds, but pushed the
duty cycle for normal teleoperation speeds below the deadband
threshold entirely, making the vehicle stop responding to real
driving commands. The gain correction was reverted
the same day to restore usable teleoperation.

\paragraph{This is a different friction phenomenon from Section~\ref{sec:friction-theory}'s -- and, in part, not friction at all.}
It is worth being precise about which effects are actually responsible
here, since this thesis initially conflated two genuinely distinct
physical phenomena with a third, unrelated software artefact, all
under similar friction-flavoured language. Section~\ref{sec:friction-theory}'s
Coulomb tyre-ground friction (the simulation's \texttt{mu1}/\texttt{mu2})
governs whether a spinning wheel's rotation converts into forward
chassis thrust at all and is essentially instantaneous. It has no
notion of a wheel being ``stuck'' and is unrelated to everything
below. The deadband characterised here is a different phenomenon
internal to the real drivetrain itself (motor brushes, gearbox mesh,
bearing seats): \emph{static} friction (often called stiction) resists
the onset of relative motion between two surfaces up to some breakaway
force threshold, while \emph{kinetic} friction, which applies once
motion has actually begun, is often measurably lower than that static
threshold for the same surface pair. A commanded duty cycle below the
static threshold produces zero motion no matter how long it is held,
since the applied torque never exceeds the force needed to break
static friction in the first place. This zero/non-zero threshold is
genuinely real and remains present after every fix in this chapter
(Section~\ref{sec:readmode-validation}).

What is \emph{not} a friction effect, despite initially being
described as one here, is the multi-second delay observed
\emph{before} that breakaway. Section~\ref{sec:readmode-bug} shows
this $\approx$\SI{2}{\second} figure is fixed regardless of duty,
speed, or how much torque is applied, something a genuine
breakaway-force process would not produce. It traces the delay
instead to a buffered sensor read in the vendor HAL.

Once duty is raised past the real static threshold, the \emph{lower}
kinetic friction then applies to an already-moving system, so the
same duty cycle that only barely broke static friction can
accelerate substantially further than expected once moving. This
asymmetry, unlike the delay, is a genuine, still-valid physical
effect. It is the stick-slip signature Section~\ref{sec:final-calibration}
identifies directly in the duty-sweep data. Table~\ref{tab:duty-sweep}'s
\texttt{0.085} row \emph{runs away} to a sustained higher
cruise speed rather than settling, a continuing behavioural
difference rather than a one-off timing offset.

This reframes the calibration problem: gain and deadband are not
independent. Any solution must jointly (a) track commanded speed
accurately at normal operating speeds and (b) keep every one of
those same commands above the static-friction threshold. A single
linear multiplier applied to the whole range cannot do both: scaling
the gain down to fix (a) simultaneously worsens (b), for exactly the
commands that were already only barely clearing the threshold.

\subsection{Final Calibration via an Empirical Duty Sweep}
\label{sec:final-calibration}

Given that finding, the calibration approach was reframed entirely:
rather than extrapolating a linear gain from one or two data points,
\texttt{THROTTLE\_GAIN} was forced to zero and a dedicated,
lighter-weight tool (\texttt{drive\_and\_log.py}) was used to sweep
fixed PWM duty values directly, with the vehicle on the ground under
real load throughout. Table~\ref{tab:duty-sweep} summarises the
resulting duty-response curve.

\begin{table}[htbp]
  \centering
  \caption{Fixed-duty sweep results (QCar on the ground, real load).}
  \label{tab:duty-sweep}
  \begin{tabular}{@{}p{2cm}p{12.5cm}@{}}
    \toprule
    Duty & Observed behaviour \\
    \midrule
    0.050 & No motion at all (confirmed twice, up to \SI{27.5}{\second} of commanding) \\
    0.060--0.061 & Reliable breakaway after an initial $\approx$\SIrange{2.0}{2.1}{\second} start-up delay (later traced to a HAL buffering artefact, Section~\ref{sec:readmode-bug}, not friction); settles to a stable, non-runaway cruise of $\approx$\SIrange{0.24}{0.33}{\meter\per\second} \\
    0.065 & $\approx$\SI{0.39}{\meter\per\second} \\
    0.085 & \textbf{Runs away} to $\approx$\SIrange{0.85}{0.9}{\meter\per\second} -- a stick-slip signature (static friction exceeds kinetic friction, so once moving, the same duty produces far more net accelerating force), confirmed by watching the full velocity trace climb continuously rather than plateau \\
    \bottomrule
  \end{tabular}
\end{table}

A caution surfaced during this sweep and worth recording explicitly:
the motion/no-motion threshold near duty $\approx 0.06$ is genuinely
noisy, not a hard deterministic cut-off. A re-test with duty
adjusted by only $+0.0013$ (\texttt{0.0533} vs.\ \texttt{0.052})
flipped a command from ``27\,s of zero motion'' to reliably covering
the full target distance. This floor appears to drift somewhat both
within and across sessions (plausibly floor contact, resting wheel
position, temperature, or battery voltage). No single threshold
measurement from this sweep should be treated as an exact, permanent
constant.

The deployed calibration combines a deadband floor with a linear gain
above it, fit from the duty=0.061/\SI{0.33}{\meter\per\second} and
duty=0.065/\SI{0.39}{\meter\per\second} points and validated
separately:

\begin{align}
  \text{duty} &= \text{THROTTLE\_DEADBAND} + \text{THROTTLE\_GAIN} \cdot |v_{\text{cmd}}|, \notag\\
  \text{THROTTLE\_DEADBAND} &= 0.04,\quad \text{THROTTLE\_GAIN} = 0.0667.
  \label{eq:throttle-calibration}
\end{align}

Listing~\ref{lst:throttle-formula} shows Equation~\ref{eq:throttle-calibration}
as implemented in \texttt{qcar\_bridge.py}'s control loop, together
with the rate limiter that caps how fast \texttt{current\_throttle}
(the value actually written to the HAL) may change per cycle. This is
a separate safeguard from the deadband/gain calibration itself,
present throughout this chapter's testing to keep duty-cycle steps
within the documented \SI{100}{\percent\per\second} hardware limit
(Appendix~\ref{app:safety-limits}).

\begin{lstlisting}[language=Python,caption={Duty-cycle command computation and rate limiting (\texttt{qcar\_bridge.py}).},label={lst:throttle-formula}]
if abs(target_linear) < 1e-3:
    desired_throttle = 0.0
else:
    desired_throttle = math.copysign(
        THROTTLE_DEADBAND + THROTTLE_GAIN * abs(target_linear), target_linear)
desired_throttle = max(-THROTTLE_LIMIT, min(THROTTLE_LIMIT, desired_throttle))
max_step = THROTTLE_RATE_LIMIT * dt
delta = max(-max_step, min(max_step, desired_throttle - current_throttle))
current_throttle += delta
\end{lstlisting}

Table~\ref{tab:calibration-validation} presents the resulting
end-to-end validation across four commanded speeds, each from a
dedicated \SI{2}{\meter} drive-and-stop test with cruise velocity
measured directly from \texttt{/odom}.

\begin{table}[htbp]
  \centering
  \caption{Throttle calibration validation (Equation~\ref{eq:throttle-calibration}).}
  \label{tab:calibration-validation}
  \begin{tabular}{@{}lccc@{}}
    \toprule
    Commanded speed & Measured cruise speed & Ratio & Note \\
    \midrule
    \SI{0.2}{\meter\per\second} & unreliable ($\approx$0.24 or no motion) & -- & At the deadband edge \\
    \SI{0.3}{\meter\per\second} & \SI{0.31}{\meter\per\second} & 1.05 & \\
    \SI{0.4}{\meter\per\second} & \SI{0.42}{\meter\per\second} & 1.05 & \\
    \SI{0.5}{\meter\per\second} & \SI{0.50}{\meter\per\second} & 1.00 & \\
    \bottomrule
  \end{tabular}
\end{table}

This is a substantial improvement over the original uncalibrated
gain's $2.6$--$2.7\times$ overshoot: commanded speeds of
\SIrange{0.3}{0.5}{\meter\per\second} now track within
$\approx$5\%. Below $\approx$\SIrange{0.25}{0.3}{\meter\per\second},
open-loop duty control against real floor/battery/temperature-dependent
friction is inherently unreliable. This is treated as a hard
physical limit of the open-loop approach at the tested operating
conditions, not a remaining tuning gap. Reliable
sub-\SI{0.25}{\meter\per\second} creep would require closed-loop
velocity feedback from the encoder, rather than a fixed duty
mapping.

\subsection{A Stray-Process False Positive}
\label{sec:stray-process}

While iterating on this calibration, a separate, purely diagnostic
false alarm is worth recording as a general debugging caution. A
prior \texttt{check\_stop\_latency.py} test process, left running in
the background rather than exiting automatically once its measurement
completed, continued publishing \texttt{linear.x=0.0} on
\texttt{/cmd\_vel} at 20\,Hz, silently racing against a separately
launched teleoperation node publishing on the same topic. The
resulting symptom was the vehicle briefly registering a throttle
command and then immediately reverting to zero within one control
cycle. It looked exactly like a gain, deadband, or logic bug and was
initially investigated as one. ROS~2 does not warn about multiple
publishers on the same topic; it simply interleaves messages, with
whichever publisher's message arrives last in a given control cycle
taking effect. The actual fix was operational, not code: kill the
stray process. This project subsequently adopted two standing
practices as a result: checking \texttt{pgrep -af} for
topic-related script names (or \texttt{ros2 topic info /cmd\_vel
--verbose} for publisher count) before deep-diving into a
control-logic theory whenever real behaviour contradicts an issued
command; and designing one-shot diagnostic nodes to exit
automatically once their measurement is complete, rather than
spinning indefinitely.

\section{The Hardware Abstraction-Layer Buffering Defect}
\label{sec:readmode-bug}

The original stop-distance symptom this chapter opened with is a
multi-second coast after a stop command, with substantial extra
distance. Even after the throttle calibration of
Section~\ref{sec:final-calibration} made commanded and real cruise
speed match closely, that symptom was still present, apparently
unaffected by the gain fix. This section documents how that residual
was ultimately traced to a defect in the vendor hardware abstraction
layer itself, not the vehicle's physical dynamics at all. This is
the project's single most significant finding, both for its
magnitude and for the methodological lesson in how long an incorrect
physical explanation can survive against evidence that, in
hindsight, directly contradicted it.

\subsection{An Initial, Incomplete Diagnosis}
\label{sec:coast-shape-analysis}

A closer look at the raw coast-phase velocity trace from one of the
calibration validation runs (\SI{0.5}{\meter\per\second}) showed the
coast was \emph{not} a smooth decay. Velocity remained essentially
flat, if anything slightly \emph{elevated} relative to the preceding
cruise average, for the first $\approx$\SI{2.0}{\second} after
``stop'' was commanded. It then crashed sharply to zero in the final
\SIrange{0.3}{0.4}{\second}. This flat-plateau-then-cliff shape
is inconsistent with passive drivetrain momentum/friction coasting
(which predicts continuous, gradual decay starting immediately once
power is cut) and instead resembles the real motor continuing to be
\emph{actively driven} for roughly two seconds after the stop command
was issued. A review of the Quanser HAL documentation
\cite{quanser2023qcar} found no
mention of any drive-motor ramp, filter, or time constant that would
explain a two-second lag of this kind -- only a separate, documented
\SI{0.16}{\second} steering-servo time constant, unrelated to the
drive motor. This shifted the leading hypothesis toward a
command-delivery lag in the project's own bridge software, rather
than physical coasting. This was a genuine change of theory,
captured here because it turned out to still be one step short of
the actual cause. A
diagnostic instrumentation pass (logging commanded velocity, the
bridge's own internally-ramped throttle value and real encoder
velocity together, all on the QCar's local clock to eliminate network
timing) was begun but interrupted by the test vehicle's battery
running out mid-session and resumed in a later session on a charged
battery.

\subsection{Disproving Two Hypotheses and the Argument That Found the Real Cause}
\label{sec:readmode-discovery}

The resumed diagnostic trace showed the bridge's own internal
throttle value reaching exactly zero within \SI{52}{\milli\second} of
the stop command, while the encoder-reported velocity stayed elevated
for roughly two more seconds. This correctly ruled out the project's
own throttle-ramping logic as the cause. As understood in hindsight,
though, it did \emph{not} actually rule out a measurement or
reporting defect in the encoder-velocity reading itself, since the
two values are produced through entirely separate code paths. This
distinction was missed on the first pass through the data.

At this point the person operating the vehicle directly and
repeatedly disputed the emerging conclusion. Teleoperating the QCar
from a genuine standstill, the vehicle visibly starts and stops moving
\emph{instantly} in response to a command. This flatly contradicts the
$\approx$2-second delay the automated test script was measuring. Two
further hypotheses were tested and disproven with direct hardware data
before the actual cause was found:

\begin{itemize}
  \item \textbf{Duty-magnitude dependence} (a higher commanded speed
    produces more torque and a faster breakaway): tested at
    \SI{0.7}{\meter\per\second} (duty 0.087, well above the normally
    tested \SIrange{0.052}{0.075} range) -- the start-up delay was
    still \SI{2.092}{\second}, statistically identical to the
    low-speed tests. Disproven.
  \item \textbf{``\texttt{v} is a stale reading, not real motion.''}
    Disproven by the structure of the \texttt{Odometry} class's own
    update code. The reported velocity \texttt{v} and the accumulated
    position \texttt{x} are computed from the exact same
    \texttt{distance} variable within the same loop iteration. If
    \texttt{v} were falsely indicating motion, the accumulated position
    would have to freeze at the same instant. It demonstrably did
    not: distance kept growing smoothly and matched the vehicle's real,
    physically measured settled position.
\end{itemize}

The argument that ultimately identified the true cause was not a
further measurement, but a piece of physical reasoning I raised
directly: \emph{an encoder is a real-time
sensor -- it cannot have a built-in two-second lag by itself.} This
was correct. It redirected the investigation to a question the
prior analysis had not actually asked: not whether the
\texttt{Odometry} class was internally self-consistent (it was), but
whether the value it was consuming as its own input, the raw
\texttt{car.motorEncoder[0]} read, was itself stale before the
\texttt{Odometry} class ever saw it. This prompted reading the vendor
HAL's own Python source directly, rather than continuing to reason
about it indirectly from documentation or from the \texttt{Odometry}
class's own behaviour.

\subsection{The Root Cause, a Ring-Buffer I/O Mode}
\label{sec:readmode-rootcause}

The root cause, found in the vendor HAL source
(\texttt{pal/products/qcar.py}) on the QCar itself, is exactly the
\texttt{readMode} distinction introduced in
Section~\ref{sec:qcar-platform}. \texttt{qcar\_bridge.py} constructed
its \texttt{QCar} object with \texttt{readMode=1} (task-based I/O).
This starts a background acquisition task \emph{at object
construction time}, before the bridge's listening TCP socket has even
opened. It fills a ring buffer of \texttt{frequency}$\times 2$
samples (at this project's \SI{50}{\hertz} read rate, exactly 100
samples, i.e.\ \SI{2.0}{\second}) with overwrite-on-overflow
behaviour. Each call to \texttt{car.read()} dequeues one buffered
sample. Because \texttt{car.read()} is not actually called until a
relay client connects, which can take many seconds after the bridge
process starts, the buffer fills and cycles during that idle gap.
Once reads begin at a consumption rate matching the production rate
exactly, the read stream is permanently stuck approximately 100
samples (\SI{2.0}{\second}) behind real time, because the backlog
never has an opportunity to drain.

\paragraph{Formalising the lag: why it is exactly the buffer depth, forever.}
Figure~\ref{fig:ringbuffer-lag} illustrates the mechanism derived
below. Let the background task write into the ring buffer at a fixed
production index $p$, incrementing by one every $1/f$ seconds
($f = \SI{50}{\hertz}$ here). Let \texttt{car.read()} calls
advance a separate consumption index $c$, each call returning the
sample at position $c \bmod B$ (buffer capacity $B = 100$) and then
incrementing $c$. Production begins at object construction, $t = 0$;
consumption begins only once a relay client connects, at
$t = T_{\text{idle}}$. During $[0, T_{\text{idle}}]$, $p$ advances
alone (production with no consumption at all), so by the time
consumption starts, $p = f \cdot T_{\text{idle}}$ while $c = 0$.
This is a gap of $f \cdot T_{\text{idle}}$ samples between them. If
$T_{\text{idle}} > B/f$ (true here, since a relay connecting takes
several seconds while $B/f = \SI{2.0}{\second}$), the overwrite-on-overflow
buffer has already wrapped around at least once by
$T_{\text{idle}}$. This means the oldest sample still available to
read is not the very first one produced, but the one produced exactly
$B$ samples before $p$'s current value. The very first
\texttt{car.read()} call, at $t = T_{\text{idle}}$, therefore
necessarily returns a sample captured at $t = T_{\text{idle}} - B/f$, a lag of
exactly $B/f = \SI{2.0}{\second}$, regardless of how much larger than
$B/f$ the idle gap $T_{\text{idle}}$ itself was. Production and
consumption both advance at the identical rate $f$ from that point
onward (one \texttt{car.read()} per control-loop cycle, matching the
background task's own sampling rate). Because of that, the gap
between $p$ and $c$ (and therefore the lag between ``now'' and
whatever sample is being read) stays fixed at exactly $B$ samples
indefinitely. Neither index ever gains on the other, so the backlog
that accumulated during the idle gap is carried forward forever
rather than draining over time. This is precisely why the lag proved
so exactly reproducible (always $B/f = \SI{2.0}{\second}$, never more
and never less, regardless of how long the QCar sat idle before a
relay connected) and why it never improved on its own the longer the
system ran.

\begin{figure}[htbp]
  \centering
  \begin{tikzpicture}[scale=1.0]
    % Axes
    \draw[-{Latex[length=2mm]}] (0,0) -- (9.5,0) node[right, font=\small] {$t$};
    \draw[-{Latex[length=2mm]}] (0,0) -- (0,5.2) node[above, font=\small] {index};

    % Production line p(t): starts at t=0, slope f, runs the whole width
    \draw[thick] (0,0) -- (9,4.8) node[right, font=\small] {production index $p$ (rate $f$)};

    % Idle gap marker on the t-axis
    \draw[dashed, gray] (3,0) -- (3,1.6);
    \node[below, font=\scriptsize] at (3,0) {$T_{\text{idle}}$};

    % Consumption line c(t): starts at t=T_idle (x=3), same slope, offset down by B (i.e. lag of B/f in time)
    \draw[thick, color=gray!70!black] (3,0) -- (9,3.2) node[right, font=\small, color=gray!70!black] {consumption index $c$ (rate $f$)};

    % Vertical gap annotation at t=7
    \draw[dotted] (7,{4.8*7/9}) -- (7,{3.2*(7-3)/6});
    \draw[{Latex[length=1.5mm]}-{Latex[length=1.5mm]}] (7.15,{4.8*7/9}) -- (7.15,{3.2*(7-3)/6})
      node[midway, right, font=\scriptsize, align=left] {$B$ samples\\$=B/f=\SI{2.0}{\second}$\\lag, constant};

    % Mark buffer wrap point
    \node[font=\scriptsize, gray, align=center] at (1.5,2.3) {buffer fills \&\\overwrites (no reader yet)};
  \end{tikzpicture}
  \caption[Ring-buffer production/consumption lag over time]{Production index $p$ (background acquisition task, running
    from $t=0$) versus consumption index $c$ (\texttt{car.read()}
    calls, starting only once a relay connects at $t=T_{\text{idle}}$).
    Both advance at the same rate $f$ once consumption begins, so the
    gap present at $t=T_{\text{idle}}$, exactly the buffer depth $B$
    once the buffer has wrapped, never closes.}
  \label{fig:ringbuffer-lag}
\end{figure}

This single mechanism explains every property of the symptom observed
across this entire chapter: the near-exact \SI{2.0}{\second} to
\SI{2.1}{\second} figure (matching the buffer's mathematically exact
2.0\,s capacity), its remarkable stability across sessions, vehicle
speeds and physical conditions (a fixed buffer depth is indifferent
to any of these) and its symmetric appearance on both the start side
(the start-up delay of Section~\ref{sec:deadband-discovery}) and the
stop side (the ``coast''), since both are read through the same
buffered path.

\textbf{Fix:} \texttt{qcar\_bridge.py} now constructs the \texttt{QCar}
object as shown in Listing~\ref{lst:readmode-fix} -- immediate I/O,
reading current hardware state directly with no background task or
buffer involved at all.

\begin{lstlisting}[language=Python,caption={The one-line HAL construction change that fixed the buffering defect (\texttt{qcar\_bridge.py}).},label={lst:readmode-fix}]
# readMode=0 (immediate I/O) - was 1 (task-based I/O), which runs a
# background acquisition task into a 100-sample ring buffer
# (frequency*2) from the moment this object is constructed,
# independent of when car.read() actually starts being called. Since
# this object is created before the listening socket even opens, and
# car.read() isn't called until a relay connects, the buffer fills and
# starts overwriting during that gap - once reads begin at the matched
# 50Hz rate, they get permanently stuck ~100 samples (~2.0s) behind
# real-time, since the backlog never drains. Immediate I/O reads
# current hardware state directly, no task/buffer involved.
car = QCar(readMode=0, frequency=READ_RATE)
\end{lstlisting}

\subsection{Validation}
\label{sec:readmode-validation}

The full calibration test suite from Section~\ref{sec:final-calibration}
was re-run in its entirety on a freshly charged battery, with the
duty formula (Equation~\ref{eq:throttle-calibration}) unchanged, to
isolate the effect of this single fix. Table~\ref{tab:readmode-fix}
presents the results.

\begin{table}[htbp]
  \centering
  \caption{Effect of the \texttt{readMode=0} fix on start-up delay, extra distance and stop latency.}
  \label{tab:readmode-fix}
  \begin{tabular}{@{}lccc@{}}
    \toprule
    Test & Start-up delay & Extra distance & Stop latency \\
    & (before $\to$ after) & (before $\to$ after) & (before $\to$ after) \\
    \midrule
    \SI{2}{\meter} @ \SI{0.3}{\meter\per\second} & \SIrange{2.09}{0.117}{\second} & \SIrange{0.520}{0.114}{\meter} & \SIrange{2.37}{0.598}{\second} \\
    \SI{2}{\meter} @ \SI{0.4}{\meter\per\second} & \SIrange{2.14}{0.001}{\second} & \SIrange{0.790}{0.114}{\meter} & \SIrange{2.40}{0.501}{\second} \\
    \SI{2}{\meter} @ \SI{0.5}{\meter\per\second} & \SIrange{2.14}{0.116}{\second} & \SIrange{0.923}{0.154}{\meter} & \SIrange{2.55}{0.550}{\second} \\
    \SI{1}{\meter} @ \SI{0.3}{\meter\per\second} & \SIrange{2.14}{0.120}{\second} & \SIrange{0.641}{0.118}{\meter} & \SIrange{2.44}{0.572}{\second} \\
    \SI{3}{\meter} @ \SI{0.3}{\meter\per\second} & \SIrange{2.08}{0.088}{\second} & \SIrange{0.583}{0.104}{\meter} & \SIrange{2.35}{0.509}{\second} \\
    \bottomrule
  \end{tabular}
\end{table}

Start-up delay collapses to \SIrange{0.001}{0.12}{\second} -- an
essentially instantaneous breakaway response, matching direct
observation of the vehicle. Extra distance becomes small and roughly
constant, $\approx$\SIrange{0.10}{0.15}{\meter}, regardless of speed
or target distance, rather than scaling with cruise speed as
previously found. This confirms that the earlier ``extra distance
scales with speed'' relationship was itself an artefact of a
fixed-\emph{duration} buffering delay applied to different real
velocities, not a genuine vehicle property.

Both the extra-distance and stop-latency figures in
Table~\ref{tab:readmode-fix} are encoder-derived, from wheel rotation
counts via \texttt{/odom}, as Section~\ref{sec:readmode-discovery}
establishes. At each of the tested distances, the encoder-reported
travel was independently cross-checked against tape-measured floor
markings and manual stopwatch timing of the commanded-stop-to-physical-stop
interval. The two agreed. This confirms the encoder as an
accurate proxy for real physical distance and timing, rather than a
further source of measurement error.

The small residual that remains after this fix is real and is fully
accounted for by ordinary physics, rather than being a further open
mystery. Dividing extra distance by stop latency, an implied average
coast velocity, yields \SIrange{54}{61}{\percent} of cruise velocity
in four of the five tests. This is consistent with an ordinary,
roughly linear velocity decay to zero over the stop-latency window (a
perfectly linear decay predicts exactly 50\%). The one outlier
(\SI{2}{\meter} @ \SI{0.3}{\meter\per\second}, 108\%) is attributable
to a data-analysis artefact (the ``cruise average'' window for that
specific test likely had not reached a true steady state, now that
the run has much less runway once start-up delay is near-instant)
rather than a physical inconsistency.

\SI{0.2}{\meter\per\second} remains genuinely unreliable even after
this fix, sometimes failing to move the vehicle at all. This is the
real, independent static-friction deadband characterised in
Section~\ref{sec:deadband-discovery} and is unrelated to the
buffering defect resolved here.

\section{Summary}
\label{sec:calibration-summary}

This chapter's investigation began with a symptom: travelling
roughly three times the commanded stop distance. This symptom had at
least three superficially plausible explanations: localisation error,
command-delivery delay and physical drivetrain coasting. It ended
by identifying two genuinely independent, compounding causes. One
was an uncalibrated, non-linear throttle-to-speed mapping
(Section~\ref{sec:throttle-calibration}). The other was a
fixed-duration sensor-read buffering artefact in the vendor hardware
abstraction layer (Section~\ref{sec:readmode-bug}) that had, for a
period of this project, been reasonably but incorrectly attributed
to genuine vehicle physics.

Neither cause was found by reasoning from documentation or first
principles alone. Both required direct, targeted measurement on the
physical robot. The second was found only after I raised a direct
physical argument that prompted reading the vendor's
own source code, rather than continuing to trust an internally
self-consistent, but incompletely verified, software abstraction.
Chapter~\ref{ch:nav-integration} builds on this
chapter's calibrated, buffering-free bridge to bring up and validate
full closed-loop Nav2 navigation on the physical vehicle.
